\section{Invariant ideals and the spectral classification}
\label{sec:classification}

We now pass from the orbit theorem to ideals of the complete deformation
ring.  Two elementary facts make this passage precise: an auxiliary formal
curve supplies a point of the required height, and translation preserves a
nonzero equation as an ambient analytic germ.  The characteristic-zero
argument uses the Gross--Hopkins period map.  We then apply the
unrestricted support comparison of Barthel--Heard--Naumann.

\Needspace{12\baselineskip}
\subsection{Points on formal subschemes}

\begin{lemma}[An auxiliary formal curve]\label{fin:curve}
Let \(R\) be a complete Noetherian local domain with residue field \(k_n\).
\begin{enumerate}
\item Suppose \(R\) is a local \(k_n\)-algebra and
\(0\ne a\in\mm_R\).  For every \(c\in\mm_C\setminus\{0\}\), there is a
continuous \(k_n\)-algebra homomorphism
\[
  R\longrightarrow\OO_C
\]
sending \(a\) to \(c\) and \(\mm_R\) into \(\mm_C\).
\item Suppose \(R\) is a local \(W(k_n)\)-algebra and \(p\ne0\) in \(R\).
Fix the usual embedding \(W(k_n)\hookrightarrow\OO_{\Cp}\).
There is a continuous \(W(k_n)\)-algebra homomorphism
\[
  R\longrightarrow\OO_{\Cp}
\]
sending \(\mm_R\) into \(\mm_{\Cp}\).
\end{enumerate}
Neither homomorphism is asserted to be injective.
\end{lemma}

\begin{proof}
We first construct a one-dimensional domain quotient on which a specified
nonzero nonunit \(a\) survives.  We use the Noetherian local dimension
and parameter theorems in \cite[Tags~00KQ and~00KD]{Stacks}.  Put \(d=\dim R\), and extend \(a\) to a
system of parameters \(a,b_2,\ldots,b_d\).  Let
\(J=(b_2,\ldots,b_d)\).  Some minimal prime \(\mathfrak q\) over \(J\)
avoids \(a\).  Indeed, otherwise \(a\in\sqrt J\), and the parameter
property would give \(\sqrt J=\mm_R\), contradicting the height bound
for an ideal generated by \(d-1\) elements.

The ring \(D=R/\mathfrak q\) is a complete local domain in which \(a\)
is a nonzero nonunit and \(\sqrt{aD}=\mm_D\).  The principal ideal
theorem gives \(\dim D\le1\), and the existence of that nonzero nonunit
gives equality.  The \(a\)-adic and maximal-ideal topologies on \(D\)
coincide.  This construction includes \(d=1\), when \(J=0\) and we may
take \(\mathfrak q=0\).

In the first case, substitution defines an injective map
\(k_n[[t]]\to D\), \(t\mapsto a\).  To check injectivity, write a nonzero
series as \(t^m v(t)\), where \(v(0)\ne0\): its image is a nonzero
power of \(a\) times a unit.  Moreover, \(D/aD\) is finite-dimensional
over \(k_n\), so completeness makes \(D\) a finite
\(k_n[[a]]\)-module.  Explicitly, lift a \(k_n\)-basis of \(D/aD\);
successively subtract linear combinations of these lifts modulo
\(a,a^2,\ldots\).  The coefficient series converge in \(k_n[[a]]\)
and express every element of \(D\) in their finite span.

For the prescribed \(c\), evaluation embeds \(k_n[[t]]\) in
\(\OO_C\): the first nonzero term of a series is uniquely dominant
when \(0<|c|<1\).  Its fraction field is a complete discretely valued
subfield of \(C\), with uniformizer \(c\).  The finite extension
\(\Frac(D)/k_n((a))\) embeds into the algebraically closed field \(C\)
over this embedding of the base.  Every element of \(D\) is integral
over \(k_n[[a]]\), so its image has norm at most one.  For
\(z\in\mm_D\), some power \(z^N\) belongs to \(aD\), whence
\[
  |z|^N\le |c|<1.
\]
Thus the composite \(R\to D\to\OO_C\) has all the required properties.

For the second case apply the same construction with \(a=p\).
The resulting \(D\) is finite over \(W(k_n)\), by the identical
successive approximation argument.  The structural map from this
DVR is injective because \(p\ne0\) in the domain \(D\).
Embed the finite extension
\(\Frac(D)/\Frac(W(k_n))\) into \(\Cp\) over the fixed base embedding.
Integrality again gives boundedness, and
\(\sqrt{pD}=\mm_D\) gives strict smallness of \(\mm_D\).

Finally, continuity follows in both cases from finite generation of
\(\mm_R\): its generators have images of norm at most some
\(\rho<1\), so \(\mm_R^N\) maps to elements of norm at most \(\rho^N\).
\end{proof}

The point supplied by Lemma~\ref{fin:curve} need not retain any given
nonzero equation after restriction to the auxiliary curve.  What matters
is that the equation retains a nonzero germ in the \emph{ambient}
polydisc.

\begin{lemma}[Bounded translation]\label{fin:translation}
Let \(L\) be a complete nontrivially valued nonarchimedean field,
let \(r\ge0\), and let
\[
  \mathscr B_r(L)=
  \left\{\sum_{I\in\mathbb N^r}a_I T^I:
                    \sup_I|a_I|<\infty\right\}.
\]
Equip this algebra with the coefficient norm
\(\|f\|=\sup_I|a_I|\).
For \(x\in L^r\) with \(\|x\|<1\), substitution \(T\mapsto x+T\)
defines an isometric \(L\)-algebra automorphism
\(\tau_x\) of \(\mathscr B_r(L)\), with inverse \(\tau_{-x}\).
In particular, a nonzero bounded-coefficient series has a nonzero
analytic germ at every such \(x\).
\end{lemma}

\begin{proof}
The cases \(r=0\) or \(x=0\) are immediate.  Otherwise put
\(\rho=\|x\|\), so \(0<\rho<1\), and
\(M=\sup_I|a_I|\).  The coefficient of \(T^J\) in the translated
series is
\[
  (\tau_x f)_J
    =\sum_{I\ge J}\binom IJ a_I x^{I-J}.
\]
Here inequalities between multi-indices are coordinatewise.
Each term has norm at most \(M\rho^{|I|-|J|}\), so the sum converges
and the new coefficients are bounded by \(M\).

For completeness, all reorderings needed for substitution can be checked
coefficientwise.  The coefficient of \(T^K\) in
\(\tau_{-x}\tau_x f\) is a sum over \(I\ge J\ge K\).
Its individual terms have norm at most
\(M\rho^{|I|-|K|}\), and only finitely many pairs occur below any
fixed total degree.  The family is therefore summable independently
of order.  Regrouping gives
\[
 \sum_{I\ge K} a_I x^{I-K}\binom IK
       \sum_{K\le J\le I}
          (-1)^{|J-K|}\binom{I-K}{J-K}
   =a_K.
\]
The inner sum is zero unless \(I=K\).  Replacing \(x\) by \(-x\)
proves the other inverse identity.

Multiplicativity has the same convergence justification.  If
\(g=\sum_Q d_QT^Q\), the coefficient of \(T^H\) in
\(\tau_x f\,\tau_x g\) is a sum over
\(J+K=H\), \(P\ge J\), and \(Q\ge K\).
The terms are bounded by
\[
  \|f\|\,\|g\|\,\rho^{|P|+|Q|-|H|},
\]
so they may be regrouped.  The multi-index Vandermonde identity
\[
  \sum_{\substack{J+K=H\\J\le P,\ K\le Q}}
           \binom PJ\binom QK=\binom{P+Q}{H}
\]
then identifies the coefficient with that of \(\tau_x(fg)\).
Both inverse maps are contractions for the coefficient supremum norm,
so they are isometries.

A bounded-coefficient series converges on every closed polydisc of
radius less than one.  Its expansion at \(x\) is the series
\(\tau_xf\), by the same convergent binomial expansion.
A zero analytic germ would have all Taylor coefficients zero and
would therefore give \(\tau_xf=0\), hence \(f=0\).
\end{proof}

\subsection{Characteristic zero and the Gross--Hopkins map}

Let \(\mathfrak X_{\Cp}\) denote the open unit polydisc with coordinates
\(u_1,\ldots,u_{n-1}\), obtained from the generic fiber of
\(\Spf(E_0)\) by extending scalars to \(\Cp\).
The nonextended stabilizer acts on it continuously by
Lemma~\ref{def:continuous-action}.  The characteristic-zero argument
below develops the Lie-theoretic approach credited to Chai in
\cite[Remark~4.2]{BHN22}; compare the projective tangent calculation in
\cite[Proposition~4.18]{BHN22}.

\begin{proposition}[Characteristic-zero orbit germs]\label{fin:charzero}
Let \(U\subseteq A_n^\times\) be an open subgroup and
\(x\in\mathfrak X_{\Cp}(\Cp)\).
If an analytic germ at \(x\) vanishes at \(g x\) for every
\(g\in U\) sufficiently close to the identity, that germ is zero.
\end{proposition}

\begin{proof}
At \(n=1\) the polydisc is a point, so the assertion is immediate.
Suppose \(n>1\).  We apply the period theorem after completed base
change from \(W(k_n)\) to \(W(k)\) and the compatible embedding
\(W(k)\to\OO_{\Cp}\).  The marked universal deformation and its
\(A_n^\times\)-action commute with this base change, and the resulting
\(\Cp\) generic fiber is \(\mathfrak X_{\Cp}\).
The Gross--Hopkins period theorem
\cite[Section~1.1, Theorem~1]{GH94}, with its proof in
\cite[Proposition~23.5]{HG94Bundles}, gives an equivariant \'etale map
\[
  \Phi:\mathfrak X_{\Cp}\longrightarrow\mathbb P^{n-1}_{\Cp}.
\]
The action on the target is the projectivization of the split
division-algebra representation
\[
  \Delta_n\otimes_{\Qp}\Cp
       \simeq\End_{\Cp}(\Cp^n).
\]
Write \(\xi=\Phi(x)\).
The analytic inverse function theorem supplies neighborhoods \(V\) of
\(x\) and \(W\) of \(\xi\) on which \(\Phi\) is an analytic isomorphism.
For this use and the one below, its local proof is the same contraction
argument used in Section~\ref{sec:orbits}: after translating and applying
the inverse linear part, write the map as \(z+R(z)\), with
\(R(0)=0\) and \(dR_0=0\), and shrink a closed polydisc until \(R\)
has Lipschitz constant less than one.
For \(w\) in that polydisc, iteration of \(z=w-R(z)\) contracts in its
complete Tate algebra, producing the analytic inverse.
Transport the given germ \(P\) to the germ
\(Q=P\circ(\Phi|_V)^{-1}\) at \(\xi\).

The infinitesimal action of the full matrix algebra surjects onto
\(T_\xi\mathbb P^{n-1}_{\Cp}\).  Since \(A_n\) spans \(\Delta_n\)
over \(\Qp\), choose \(d_1,\ldots,d_{n-1}\in A_n\) whose infinitesimal
projective directions form a \(\Cp\)-basis of this tangent space.
For small \(t=(t_1,\ldots,t_{n-1})\in\Cp^{n-1}\), form
\[
  \eta(t)=\left(1+\sum_i t_i d_i\right)\xi,
\]
using the split representation to interpret the matrix on the right.
This is an analytic map with invertible differential at zero.
Shrink its domain so that its image lies in \(W\) and \(Q\circ\eta\)
is defined.

For sufficiently large \(a\), every
\(t\in(p^a\Zp)^{n-1}\) defines an element
\(g_t=1+\sum_i t_i d_i\in U\).
Continuity of the stabilizer action allows us also to require
\(g_t x\in V\), within the neighborhood on which \(P\) vanishes along
the orbit.  Equivariance and the fixed inverse branch give
\[
  (\Phi|_V)^{-1}(\eta(t))=g_t x,
  \qquad (Q\circ\eta)(t)=0.
\]
An analytic function on a \(\Cp\)-polydisc that vanishes on
\((p^a\Zp)^{n-1}\) is zero: fix all but one variable in that set,
apply one-variable uniqueness using a sequence tending to zero,
and repeat for each variable.  Indeed a nonzero one-variable germ has
the form \(t^r u(t)\) with \(u(0)\ne0\), so its zeros cannot
accumulate at zero.  Thus \(Q\circ\eta\) is the zero germ.
The invertible differential of \(\eta\) implies that \(Q\), and
therefore \(P\), is the zero germ.
\end{proof}

This argument uses a \(\Qp\)-valued set of parameters in a fixed
analytic inverse branch.  It requires no openness of the stabilizer
orbit in the \(\Cp\)-topology.

\subsection{Primes and radical ideals}

\begin{theorem}\label{fin:primes}
Every proper prime ideal of \(E_0\) stable under an open subgroup of
\(A_n^\times\) is one of
\[
  0,\ I_1,\ldots,I_n.
\]
\end{theorem}

\begin{proof}
Let \(\mathfrak p\) be such a prime, stable under the open subgroup
\(U\).

Suppose first that \(p\in\mathfrak p\), and choose the largest \(h\)
such that \(I_h\subseteq\mathfrak p\).  If \(h=n\), the ideal \(I_n\)
is maximal, so \(\mathfrak p=I_n\).  Otherwise the image of \(u_h\)
in \(R=E_0/\mathfrak p\) is a nonzero nonunit.
Lemma~\ref{fin:curve} gives a point \(x:R\to\OO_C\) at which all
parameters are small and \(u_h(x)\ne0\).
It is therefore a point of the height-at-least-\(h\) polydisc to
which Theorem~\ref{orb:density} applies.

If \(\mathfrak p/I_h\ne0\), choose a nonzero series
\[
  f\in\mathfrak p/I_h
     \subset k_n[[u_h,\ldots,u_{n-1}]].
\]
Its coefficients are bounded, and Lemma~\ref{fin:translation} shows
that it defines a nonzero ambient analytic germ at \(x\).
Stability of \(\mathfrak p\) makes \(f\) vanish at every point of the
\(U\)-orbit of \(x\).
This contradicts Theorem~\ref{orb:density}.
Consequently \(\mathfrak p=I_h\).

Now suppose \(p\notin\mathfrak p\).
The mixed-characteristic part of Lemma~\ref{fin:curve} supplies a point
\(x:E_0/\mathfrak p\to\OO_{\Cp}\) with all parameters small.
If \(\mathfrak p\ne0\), take \(0\ne f\in\mathfrak p\).
The fixed embedding \(W(k_n)\hookrightarrow\OO_{\Cp}\) regards \(f\)
as a nonzero bounded-coefficient series in the ambient generic
polydisc.  By Lemma~\ref{fin:translation}, its germ at \(x\) is
nonzero.  It vanishes on the \(U\)-orbit, contradicting
Proposition~\ref{fin:charzero}.  Thus \(\mathfrak p=0\).
\end{proof}

\begin{corollary}\label{fin:radicals}
For every open subgroup \(U\subseteq A_n^\times\), the
\(U\)-invariant radical ideals of \(E_0\) are precisely
\[
  0,\ I_1,\ldots,I_n,\ E_0.
\]
The same list gives the radical ideals invariant under the full
canonical extended stabilizer group \(G_n\).
\end{corollary}

\begin{proof}
Let \(J\) be a proper \(U\)-invariant radical ideal.
Noetherianity gives finitely many minimal primes, and the group
permutes them.  The stabilizer of each such prime \(P\) is closed.
Indeed, ideals in the complete Noetherian local ring \(E_0\) are
adically closed, and each of the conditions
\[
  g(a)\in P,\qquad g^{-1}(a)\in P \qquad(a\in P)
\]
is closed by continuity of the action.  Together these conditions
say \(gP=P\).  This closed stabilizer has finite index in \(U\) and is
therefore open in \(U\): its complement is a finite union of closed
cosets.  Since \(U\) is open in \(A_n^\times\), the stabilizer is
open in \(A_n^\times\) as well.
Theorem~\ref{fin:primes} applies to every minimal prime of \(J\).
Their intersection is \(J\); since the standard primes form a chain,
this intersection is its smallest member.  Hence \(J\) is standard.

An ideal invariant under \(G_n\) is invariant under
\(A_n^\times\).  Conversely, every ideal in the displayed list is
invariant under the full canonical action.  For \(I_h\), this follows
from its intrinsic description as the locus where the universal
formal group has height at least \(h\), or equivalently from the
successive height parameters of its \(p\)-series.  The zero and unit
ideals are invariant as well.
\end{proof}

Theorem~\ref{fin:primes} and Corollary~\ref{fin:radicals} prove
Theorem~\ref{thm:invariant-radicals}.

\subsection{Support and the spectral classification}
\label{fin:classification}

Write \(\mathcal D=\mathcal D_{(n,p)}\), and denote the thick tensor
ideal generated by an object \(X\) by \(\langle X\rangle_\otimes\).
For \(X\in\mathcal D\), put
\[
  \mathcal F(X)=L_{K(n)}(E\wedge X),\qquad
  M_X=\pi_0\mathcal F(X)\oplus\pi_1\mathcal F(X),\qquad
  S(X)=\Supp_{E_0}(M_X).
\]
The two summands of \(M_X\) are the completed Morava \(E\)-homology
groups in degrees zero and one, which determine the remaining degrees
by periodicity.

We recall precisely the support input needed here.
It is the all-prime direction of the reduction in
\cite[Theorem~4.9]{BHN22}.

\begin{proposition}[Established support comparison]\label{fin:support}
For every prime \(p\) and height \(n\ge1\), the following hold.
\begin{enumerate}
\item The functor
\(\mathcal F:\mathcal D\to\Perf(E)\) is exact and symmetric monoidal,
where the source uses the \(K(n)\)-localized product.
It detects tensor-nilpotence of morphisms.
\item The set \(S(X)\) is a closed \(G_n\)-invariant subset
of \(\Spec(E_0)\), and
\[
 S(X\otimes Y)=S(X)\cap S(Y).
\]
For \(X,Y\in\mathcal D\),
\[
  X\in\langle Y\rangle_\otimes
       \quad\Longleftrightarrow\quad
  S(X)\subseteq S(Y).
\]
\item For any finite \(p\)-local spectrum \(F(k)\) of type \(k\),
\(0\le k\le n\),
\[
  S\bigl(L_{K(n)}F(k)\bigr)=V(I_k).
\]
Also \(S(0)=\varnothing\) and
\(S(L_{K(n)}S)=\Spec(E_0)\).
\end{enumerate}
\end{proposition}

\begin{proof}
Here we assemble the established results and spell out the direction
of support descent.
Base change is symmetric monoidal into the category of
\(K(n)\)-local \(E\)-modules.  Its dualizable objects are exactly
the ordinary perfect \(E\)-modules by
\cite[Proposition~4.7]{BHN22}.  Thus \(\mathcal F\) has the asserted
target even when \(X\) is dualizable but not compact as a
\(K(n)\)-local spectrum.  Tensor-nilpotence detection is
\cite[Proposition~4.6]{BHN22}; neither statement has a prime bound.
These inputs use Morava-\(E\) descent and the perfect-module comparison
of Mathew \cite[Propositions~3.27 and~10.10--10.11]{Mathew16}.

The coefficient ring \(E_*\) is even-periodic, regular, and
Noetherian.  The comparison theorem of Dell'Ambrogio--Stanley
\cite[Theorem~1.1 and Lemma~3.10]{DellAmbrogioStanley16}, applied to
Morava \(E\)-theory in \cite[Lemmas~4.3--4.4]{BHN22}, identifies
\[
  \Spc(\Perf(E))\simeq\Spec(E_0)
\]
and identifies the tensor-triangular support of \(\mathcal F(X)\)
with \(S(X)\).  The tensor-product identity for that support gives
\(S(X\otimes Y)=S(X)\cap S(Y)\).
The module \(M_X\) is finite over \(E_0\), so this
support is closed.  Its canonical semilinear stabilizer action makes
it invariant; equivalently one may use
\cite[Lemma~4.5]{BHN22}.

Tensor-nilpotence detection gives a surjection
\[
 f:\Spc(\Perf(E))\longrightarrow\Spc(\mathcal D)
\]
by Balmer's surjectivity theorem \cite[Theorem~1.3]{Balmer18},
in the form \cite[Theorem~2.6 and Corollary~4.8]{BHN22}.
Writing \(\operatorname{supp}\) for tensor-triangular support,
functoriality gives
\[
 \operatorname{supp}(\mathcal F(X))
       =f^{-1}\bigl(\operatorname{supp}(X)\bigr).
\]
If \(S(X)\subseteq S(Y)\), surjectivity of \(f\) therefore implies
\(\operatorname{supp}(X)\subseteq\operatorname{supp}(Y)\).
Every object of \(\mathcal D\) is dualizable, so its thick tensor
ideals are radical.  Balmer's classification
\cite[Theorem~4.10]{Balmer05}, in the support-containment form of
\cite[Corollary~2.4]{BHN22}, gives
\(X\in\langle Y\rangle_\otimes\).
The converse follows from exactness, tensor compatibility, and the
same support criterion.  This is also the transfer formulated in
\cite[Corollary~2.7]{BHN22}.

For the finite type-\(k\) objects the coefficient support is
\(V(I_k)\), as in
\cite[Proposition~3.3 and the proof of Theorem~4.9]{BHN22}.
This is the usual finite chromatic support calculation; it is
unrestricted in \(p\).
The two remaining support calculations follow from
\(\mathcal F(0)=0\) and \(\mathcal F(L_{K(n)}S)=E\).
\end{proof}

\begin{proof}[Proof of Theorem~\ref{thm:main}]
By Corollary~\ref{fin:radicals} and
Proposition~\ref{fin:support}, the support of every object of
\(\mathcal D\) belongs to the finite chain
\[
 \Spec(E_0)=V(I_0)\supsetneq V(I_1)\supsetneq\cdots
       \supsetneq V(I_n)\supsetneq\varnothing.
\]
Indeed, \(S(X)=V(\Ann_{E_0}M_X)\), and its defining radical ideal is
\(G_n\)-invariant.
Empty support occurs only for the zero object, by applying the
support comparison with \(Y=0\).

Let \(\mathcal J\ne0\) be any thick tensor ideal.
The supports occurring among its objects form a nonempty subset
of this finite chain, and so have a greatest member that is
\emph{attained} by some \(Y\in\mathcal J\).
For every \(X\in\mathcal J\), support comparison gives
\(X\in\langle Y\rangle_\otimes\).
The reverse containment follows from \(Y\in\mathcal J\), hence
\[
  \mathcal J=\langle Y\rangle_\otimes.
\]
There is no assumption that \(\mathcal J\) was singly generated:
attainment follows solely from finiteness of the possible supports.

Write \(S(Y)=V(I_k)\).
The given object \(L_{K(n)}F(k)\) has the same support, so applying
Proposition~\ref{fin:support} in both directions gives
\[
  \mathcal J
    =\langle L_{K(n)}F(k)\rangle_\otimes
    =\mathcal D_k.
\]
The zero ideal is \(\mathcal D_{n+1}\).

The same comparison proves independence of the generator:
any other finite type-\(k\) spectrum has support \(V(I_k)\) and
therefore generates the same ideal.
This also agrees with the finite thick-subcategory theorem of
Hopkins--Smith \cite{HopkinsSmith98}.
The strict inclusions of the displayed supports imply
\(\mathcal D_k\supsetneq\mathcal D_{k+1}\), since
\(L_{K(n)}F(k)\) belongs to the first ideal and its support is not
contained in that of a generator of the second.
For \(k=n\), its support \(V(I_n)\) is nonempty, so
\(\mathcal D_n\ne0\).
Finally \(V(I_0)=\Spec(E_0)\) is the support of the tensor unit.
Thus \(\mathcal D_0\) contains the unit and equals \(\mathcal D\).
The list has exactly \(n+2\) distinct entries, as asserted.
\end{proof}

Every support needed in this deduction is already realized by the
finite spectra specified in the theorem.  No realization of an
arbitrary Morava module, and hence no prime restriction from the
converse direction of the Hovey--Strickland--Chai comparison, enters
the argument.


\subsection{The Balmer spectrum}

A proper thick tensor ideal $\mathcal P$ is \emph{prime} if
$X\otimes Y\in\mathcal P$ implies $X\in\mathcal P$ or
$Y\in\mathcal P$.  The points of $\Spc(\mathcal D)$ are these ideals;
the tensor-triangular support of $X$ is
$\operatorname{supp}(X)=\{\mathcal P:X\notin\mathcal P\}$.
The following identifies the spectrum, with the indexing used in
\cite[Proposition~3.5]{BHN22}.

\begin{corollary}\label{fin:balmer}
For every prime $p$ and every $n\geq1$,
\[
 \Spc(\mathcal D_{(n,p)})
      =\{\mathcal D_1,\ldots,\mathcal D_{n+1}\},\qquad
 \overline{\{\mathcal D_j\}}
      =\{\mathcal D_l:j\leq l\leq n+1\}.
\]
Thus $\mathcal D_1$ is the generic point and
$\mathcal D_{n+1}=0$ is the unique closed point.  For $0\leq k\leq n$,
\[
 \operatorname{supp}(L_{K(n)}F(k))
      =\{\mathcal D_j:k<j\leq n+1\}.
\]
\end{corollary}

\begin{proof}
The possible coefficient supports form a total chain, including the
empty set.  Hence $S(X)\cap S(Y)$ is one of $S(X)$ and $S(Y)$.
By the tensor-product and containment statements in
Proposition~\ref{fin:support}, if $X\otimes Y$ belongs to a proper
$\mathcal D_j$, then one of $X$ and $Y$ does as well.
This also includes $\mathcal D_{n+1}=0$: two nonempty supports both
contain $V(I_n)$, so their intersection is nonempty.
Every proper member of the chain is therefore prime, and
Theorem~\ref{thm:main} leaves no other prime ideals.

In the Balmer topology, a point $\mathcal Q$ lies in
$\overline{\{\mathcal P\}}$ precisely when
$\mathcal Q\subseteq\mathcal P$.
Indeed every support containing $\mathcal P$ must contain
$\mathcal Q$, which says that $X\notin\mathcal P$ implies
$X\notin\mathcal Q$ for every object $X$.
The strict ideal chain gives the stated closure formula.
Finally $L_{K(n)}F(k)\in\mathcal D_j$ exactly when $k\geq j$,
by the coefficient-support comparison, proving the support formula.
\end{proof}

The continuous surjection in Proposition~\ref{fin:support} has an
explicit description in terms of the height ideals.  For
$\mathfrak q\in\Spec(E_0)$, set
\[
 h(\mathfrak q)=\max\{0\leq h\leq n:I_h\subseteq\mathfrak q\}.
\]
Then
\[
 f:\Spec(E_0)\longrightarrow\Spc(\mathcal D),\qquad
 f(\mathfrak q)=\mathcal D_{h(\mathfrak q)+1}.
\]
Indeed, functoriality of support says that
$X\in f(\mathfrak q)$ exactly when $\mathfrak q\notin S(X)$.
For a nonzero object with $S(X)=V(I_k)$, this is equivalent to
$k>h(\mathfrak q)$, which by Proposition~\ref{fin:support} is
exactly the condition $X\in\mathcal D_{h(\mathfrak q)+1}$.
The zero object belongs to both ideals as well.
Consequently the fiber over $\mathcal D_{h+1}$ is
$V(I_h)\setminus V(I_{h+1})$ for $0\leq h<n$, and the fiber over
$\mathcal D_{n+1}=0$ is $V(I_n)$.
Here $h(\mathfrak q)$ records the formal-group height stratum, not the
Krull height of the prime: all primes avoiding $p$ map to the generic
Balmer point $\mathcal D_1$, while the maximal ideal $I_n$ maps to
the closed point $0$.
